\chapter{Shing-Tung Yau --- Fields Medal 1982}\label{ch1982shingtungyau:shing-tung-yau-fields-medal-1982-a-comprehensive-report-on-his-work-and-impact}

\begin{center}\rule{0.5\linewidth}{0.5pt}\end{center}

\textbf{Prepared:} July 2026 \textbf{Subject:} Shing-Tung Yau (b. April
4, 1949, Shantou, Guangdong, China) \textbf{Position:} Professor,
Harvard University; Tsinghua University \textbf{Doctoral Advisor:}
Shiing-Shen Chern (UC Berkeley) \textbf{Fields Medal Citation:} ``For
his contributions to differential equations, to the Calabi conjecture in
algebraic geometry, to the positive mass conjecture of general
relativity theory, and to real and complex Monge--Ampère equations.''

\begin{center}\rule{0.5\linewidth}{0.5pt}\end{center}

\section{Table of Contents}\label{ch1982shingtungyau:table-of-contents}

\begin{enumerate}
\def\labelenumi{\arabic{enumi}.}
\item
 \hyperref[introduction]{Introduction}
\item
 \hyperref[biographical-background]{Biographical Background}
\item
 \hyperref[geometric-analysis-overview]{Geometric Analysis: Overview}
\item
 \hyperref[the-calabi-conjecture-background]{The Calabi Conjecture:
 Background}
\item
 \hyperref[the-calabi-conjecture-proof-and-solutions]{The Calabi
 Conjecture: Proof and Solutions}
\item
 \hyperref[calabi-yau-manifolds]{Calabi-Yau Manifolds}
\item
 \hyperref[the-positive-mass-conjecture-background]{The Positive Mass
 Conjecture: Background}
\item
 \hyperref[the-positive-mass-conjecture-minimal-surface-proof]{The
 Positive Mass Conjecture: Minimal Surface Proof}
\item
 \hyperref[real-monge-ampere-equations]{Real Monge-Ampere Equations}
\item
 \hyperref[complex-monge-ampere-equations]{Complex Monge-Ampere
 Equations}
\item
 \hyperref[the-yamabe-problem-and-kahler-einstein-metrics]{The Yamabe
 Problem and Kahler-Einstein Metrics}
\item
 \hyperref[the-smith-conjecture]{The Smith Conjecture}
\item
 \hyperref[differential-geometry-of-submanifolds]{Differential Geometry
 of Submanifolds}
\item
 \hyperref[black-hole-formation-and-general-relativity]{Black Hole
 Formation and General Relativity}
\item
 \hyperref[mirror-symmetry-and-string-theory]{Mirror Symmetry and
 String Theory}
\item
 \hyperref[unifying-techniques-in-geometric-analysis]{Unifying
 Techniques in Geometric Analysis}
\item
 \hyperref[broad-impact-on-modern-mathematics]{Broad Impact on Modern
 Mathematics}
\item
 \hyperref[awards-and-professional-recognition]{Awards and Professional
 Recognition}
\item
 \hyperref[selected-major-publications]{Selected Major Publications}
\item
 \hyperref[ongoing-legacy-and-future-directions]{Ongoing Legacy and
 Future Directions}
\item
 \hyperref[conclusion]{Conclusion}
\item ~
 \section{\texorpdfstring{\hyperref[references-and-further-reading]{References
 and Further
 Reading}}{References and Further Reading}}\label{ch1982shingtungyau:references-and-further-reading}
\end{enumerate}

\begin{prereqs} Differential geometry (Riemannian manifolds, curvature), basic general relativity, partial differential equations (elliptic theory). \end{prereqs}

\begin{objectives} Explain the Calabi conjecture and its proof; describe the positive mass theorem; understand the role of geometric analysis in physics. \end{objectives}


\section{1. Introduction}\label{ch1982shingtungyau:introduction}

\subsection{1.1 Overview of Legacy}\label{ch1982shingtungyau:overview-of-legacy}

Shing-Tung Yau is widely recognized as one of the influential
mathematicians of modern times. He is the principal founder of geometric
analysis, a discipline that integrates differential geometry, complex
geometry, topology, and the theory of nonlinear partial differential
equations (PDEs). Over a career spanning more than five decades, Yau has
solved several of the well-known problems in mathematics and
theoretical physics.

His work has established a significant connection between the abstract
structure of mathematical manifolds and the physical equations
describing gravity and spacetime. By treating geometric questions
through the lens of nonlinear analysis, Yau unlocked new methods to
study the shape and stability of the universe.

\subsection{1.2 The Fields Medal of
1982}\label{ch1982shingtungyau:the-fields-medal-of-1982}

In 1982, Yau was awarded the Fields Medal, the highest honor in
mathematics, for his pioneering contributions. Although the
International Congress of Mathematicians (ICM) in Warsaw was delayed to
1983 due to martial law in Poland, Yau is officially recorded as the
1982 medalist.

The citation for his medal specifically noted his contributions to: -
The Calabi conjecture in algebraic geometry. - The positive mass
conjecture in general relativity. - Real and complex Monge--Ampère
equations. - The theory of nonlinear differential equations.

As the first mathematician of Chinese descent to win the Fields Medal,
his achievement served as an inspiration for mathematicians worldwide,
paving the way for the growth of geometric analysis across many research
institutions.

\subsection{1.3 Bridging Mathematics and
Physics}\label{ch1982shingtungyau:bridging-mathematics-and-physics}

Yau's research is characterized by a unique combination of analytic
power and geometric intuition. He has maintained a productive dialogue
with theoretical physicists, demonstrating that mathematical rigor and
physical intuition can enrich one another.

His proof of the Calabi conjecture provided string theorists with the
geometric building blocks---Calabi-Yau manifolds---required for the
compactification of extra dimensions. Meanwhile, his proof of the
positive mass conjecture established the mathematical stability of
Einstein's general theory of relativity. Yau's legacy is defined by this
capability to translate physical questions into rigorous geometry, and
vice versa.

\subsection{1.4 Structure of the Report}\label{ch1982shingtungyau:structure-of-the-report}

This report is structured as a comprehensive survey of Shing-Tung Yau's
contributions to mathematics and theoretical physics. Sections 2 and 3
establish his biographical and methodological foundations, while
Sections 4 through 6 detail the Calabi conjecture and the manifolds
named after him.

Sections 7 and 8 cover the positive mass conjecture and the minimal
surface techniques used in its proof. Sections 9 and 10 discuss his
contributions to Monge-Ampère equations, leading to the analysis of the
Yamabe problem and Kähler-Einstein metrics in Section 11.

Sections 12 and 13 cover his work on the Smith conjecture and
submanifold geometry, while Sections 14 and 15 detail his contributions
to mathematical physics, including black holes and mirror symmetry.
Sections 16 and 17 evaluate his unifying mathematical techniques and
broad institutional impact.

Finally, Sections 18 and 19 list his awards and major publications,
followed by an overview of his ongoing legacy in Section 20, a general
conclusion in Section 21, and references in Section 22.

\begin{center}\rule{0.5\linewidth}{0.5pt}\end{center}

\section{2. Biographical Background}\label{ch1982shingtungyau:biographical-background}

\subsection{2.1 Early Life in China and Hong
Kong}\label{ch1982shingtungyau:early-life-in-china-and-hong-kong}

Shing-Tung Yau was born on April 4, 1949, in Shantou, Guangdong, China.
Shortly after his birth, his family relocated to Hong Kong amidst the
political instability of the Chinese Civil War. Growing up in the rural
areas of Hong Kong, including Yuen Long and Sha Tin, Yau experienced
severe financial hardship.

His family was large, consisting of eight children. His father,
Chen-Ying Yau, was a professor of classical literature and philosophy,
whose intellectual influence closely shaped Yau's early life. His father
taught him classical Chinese poetry, history, and western philosophy,
instilling in him a love for scholarly inquiry and deep reflection.

\subsection{2.2 The Passing of Chen-Ying
Yau}\label{ch1982shingtungyau:the-passing-of-chen-ying-yau}

In 1963, when Yau was only fourteen years old, his father passed away.
This tragedy plunged the family into deep poverty. To help support his
mother and siblings, Yau worked as a tutor while continuing his own
secondary education.

Despite these hardships, he excelled academically. He attended Pui Ching
Middle School, where his mathematical talent was recognized by his
teachers. They provided him with advanced mathematics books, allowing
him to study calculus and classical geometry independently, far ahead of
his school peers.

\subsection{2.3 Higher Education in Hong Kong and
Berkeley}\label{ch1982shingtungyau:higher-education-in-hong-kong-and-berkeley}

In 1966, Yau entered the Chinese University of Hong Kong (CUHK),
majoring in mathematics. He completed his undergraduate coursework in
three years. During this time, he met the famous geometer Shiing-Shen
Chern, who was visiting Hong Kong.

Chern recognized Yau's exceptional mathematical talent and secured a
scholarship for him to pursue graduate studies at the University of
California, Berkeley. Yau arrived at Berkeley in the fall of 1969. Under
Chern's supervision, Yau completed his PhD in 1971 at the age of
twenty-two. His doctoral thesis, titled ``Subgroups of the fundamental
group of a manifold of nonpositive curvature,'' demonstrated his early
mastery of the intersection between geometry and topology.

\subsection{2.4 Academic Career and Key
Appointments}\label{ch1982shingtungyau:academic-career-and-key-appointments}

Following his PhD, Yau spent a year as a member of the Institute for
Advanced Study (IAS) in Princeton (1971--1972). He then joined the
mathematics department at the State University of New York at Stony
Brook as an assistant professor (1972--1973).

In 1974, he moved to Stanford University, where he was promoted to full
professor in 1978. In 1979, he returned to the Institute for Advanced
Study as a professor. In 1984, he transitioned to the University of
California, San Diego.

In 1987, Yau joined the faculty of Harvard University, where he served
as the William Caspar Graustein Professor of Mathematics until his
retirement in 2022. Following his retirement from Harvard, Yau accepted
a full-time position at Tsinghua University in Beijing as the Director
of the Yau Mathematical Sciences Center (YMSC), where he continues to
lead mathematical research.

\subsection{2.5 Intellectual Heritage and Academic
Lineage}\label{ch1982shingtungyau:intellectual-heritage-and-academic-lineage}

Through his advisor, Shiing-Shen Chern, Yau inherited a rich
mathematical lineage. Chern was a student of Wilhelm Blaschke and Élie
Cartan, two of the most prominent geometers of the early twentieth
century.

This lineage connected Yau to classical differential geometry, which
focused on local properties of curves and surfaces, and Cartan's theory
of moving frames. Yau took these classical ideas and combined them with
the modern theory of partial differential equations, creating a new
approach that transformed differential geometry into a global, analytic
discipline.

\begin{center}\rule{0.5\linewidth}{0.5pt}\end{center}

\section{3. Geometric Analysis:
Overview}\label{ch1982shingtungyau:geometric-analysis-overview}

\subsection{3.1 Unification of Analysis and
Geometry}\label{ch1982shingtungyau:unification-of-analysis-and-geometry}

Prior to the 1970s, differential geometry and partial differential
equations were studied as separate disciplines. Geometers relied
primarily on algebraic topology, complex analysis, and linear
invariants, while analysts studied elliptic and parabolic PDEs in
Euclidean domains.

Yau, alongside collaborators like Richard Schoen, Leon Simon, and Karen
Uhlenbeck, pioneered the integration of these fields, establishing
geometric analysis. The core philosophy of geometric analysis is that
the topological and geometric properties of a manifold are determined by
the behavior of physical fields or geometric equations defined on them.

\subsection{3.2 Role of Nonlinear Partial Differential
Equations}\label{ch1982shingtungyau:role-of-nonlinear-partial-differential-equations}

In geometric analysis, geometric invariants like curvature, volume, and
area are represented as coefficients and variables in differential
equations. For instance, the problem of finding a metric with constant
curvature leads to a nonlinear elliptic PDE.

By studying these equations, mathematicians can determine whether a
manifold can support specific geometric structures. Yau demonstrated
that the analytic properties of PDEs (such as solvability, regularity,
and boundary behavior) are closely connected to the topology and geometry
of the underlying space.

\subsection{3.3 The Core Toolkit: A Priori Estimates and Maximum
Principles}\label{ch1982shingtungyau:the-core-toolkit-a-priori-estimates-and-maximum-principles}

The primary challenge in geometric analysis is that the governing PDEs
are almost always fully nonlinear. To prove that a solution exists, one
must establish a priori estimates---bounds on the solution and its
derivatives that prevent it from blowing up or degenerating.

Yau developed a suite of analytical tools to achieve this: -
\textbf{Yau's Maximum Principle:} An extension of the classical Hopf
maximum principle to complete, non-compact Riemannian manifolds with
curvature bounded from below. - \textbf{Moser Iteration:} An
algebraic-analytic technique used to bootstrap weak \(L^p\) bounds of a
solution into strong \(L^\infty\) bounds. - \textbf{Gradient Estimates:}
Local and global bounds on the derivative of a solution, which are
crucial for ensuring the regularity and smoothness of the metric or
surface under study.

\subsection{3.4 Historical Precedents and the Transition from Classical
Methods}\label{ch1982shingtungyau:historical-precedents-and-the-transition-from-classical-methods}

Before the advent of geometric analysis, the main tool for relating
curvature to topology was the Gauss-Bonnet theorem and its
generalizations (like the Chern-Weil theory). These methods were
topological and algebraic in nature, relying on integration of curvature
forms.

While powerful, they could not construct new metrics or determine the
existence of special shapes (like minimal surfaces) on arbitrary
manifolds. The transition to geometric analysis allowed mathematicians
to deform metrics dynamically and solve boundary value problems,
providing a constructive approach to geometry that did not exist before.

\begin{center}\rule{0.5\linewidth}{0.5pt}\end{center}

\section{4. The Calabi Conjecture:
Background}\label{ch1982shingtungyau:the-calabi-conjecture-background}

\subsection{4.1 Kähler Manifolds and
Curvature}\label{ch1982shingtungyau:kuxe4hler-manifolds-and-curvature}

To understand the Calabi conjecture, one must first define the setting
of Kähler geometry. A Kähler manifold \(M\) is a complex manifold
equipped with a Riemannian metric \(g\) that is compatible with the
complex structure \(J\), meaning: \[ g(JX, JY) = g(X, Y) \] The
compatibility allows one to define a closed differential 2-form called
the Kähler form: \[ \omega(X, Y) = g(JX, Y), \quad d\omega = 0 \] On a
Kähler manifold, the Ricci tensor represents the curvature of the
canonical line bundle. The Ricci form \(\text{Ric}(\omega)\), defined in
local coordinates by:
\[ \text{Ric}(\omega) = -i \partial \bar{\partial} \log \det(g_{a\bar{b}}) \]
represents the first Chern class \(c_1(M)\) of the manifold, which is a
topological invariant.

\subsection{4.2 Eugenio Calabi's
Conjecture}\label{ch1982shingtungyau:eugenio-calabis-conjecture}

In 1953 and 1954, Eugenio Calabi proposed a conjecture concerning the
relationship between the first Chern class and the Ricci curvature of
compact Kähler manifolds. Calabi asked whether the topological class
\(c_1(M)\) is the only obstruction to prescribing the Ricci curvature of
a Kähler manifold.

Formally, let \(M\) be a compact Kähler manifold with Kähler form
\(\omega\). If \(\Omega\) is any closed real \((1,1)\)-form belonging to
the cohomology class \(2\pi c_1(M)\), Calabi conjectured that there
exists a unique Kähler form \(\tilde{\omega}\) in the same cohomology
class as \(\omega\) (meaning \([\tilde{\omega}] = [\omega]\)) such that
its Ricci form is precisely \(\Omega\):
\[ \text{Ric}(\tilde{\omega}) = \Omega \] If this conjecture were true,
it would mean that we could prescribe the Ricci curvature of a compact
Kähler manifold almost at will, subject only to the topological
constraint of the first Chern class.

\subsection{4.3 Yau's Journey: From Counterexamples to
Proof}\label{ch1982shingtungyau:yaus-journey-from-counterexamples-to-proof}

For over two decades, the Calabi conjecture remained one of the most
famous unsolved problems in complex geometry. Many geometers, including
Yau in his early years, believed the conjecture was false and tried to
construct counterexamples.

In 1973, at a conference at Stanford University, Yau announced that he
had found a counterexample. However, Eugenio Calabi wrote to Yau,
pointing out a gap in the proof. In attempting to fix the gap, Yau
realized that all his candidate counterexamples actually satisfied the
conjecture. This insight convinced him that Calabi's conjecture must be
true, and he spent the next several years developing the analytical
machinery required to prove it.

\subsection{4.4 The 1954 ICM Address in
Amsterdam}\label{ch1982shingtungyau:the-1954-icm-address-in-amsterdam}

Eugenio Calabi formally presented his conjecture at the International
Congress of Mathematicians in Amsterdam in 1954. In his address, he
outlined the mathematical consequences of the conjecture, showing that
it would imply the existence of metrics with vanishing Ricci curvature
on manifolds with zero first Chern class.

Despite Calabi's clear formulation, the mathematical community remained
skeptical because there were no general methods for solving nonlinear
elliptic equations on manifolds at the time. The conjecture was viewed
as a beautiful but inaccessible dream until Yau's breakthrough
twenty-four years later.

\subsection{4.5 Skepticism of the Mathematical
Community}\label{ch1982shingtungyau:skepticism-of-the-mathematical-community}

In the decades following Calabi's address, several attempts were made to
find counterexamples on specific algebraic surfaces. Leading
mathematicians of the era, such as René Thom and Harish-Chandra,
expressed doubts about the global existence of Ricci-flat metrics.

The main argument against the conjecture was that a Ricci-flat metric on
a compact manifold without boundaries seemed too rigid to exist. Yau's
work dispelled these doubts, showing that the geometric rigidity was
matched by a corresponding analytical flexibility in the underlying
complex Monge-Ampère equation.

\begin{center}\rule{0.5\linewidth}{0.5pt}\end{center}

\section{5. The Calabi Conjecture: Proof and
Solutions}\label{ch1982shingtungyau:the-calabi-conjecture-proof-and-solutions}

\subsection{5.1 Formulation as a Complex Monge-Ampère
Equation}\label{ch1982shingtungyau:formulation-as-a-complex-monge-ampuxe8re-equation}

Yau's proof of the Calabi conjecture, published in 1978, relies on
translating the geometric problem into a fully nonlinear complex
Monge-Ampère PDE. Let \(\omega\) be the initial Kähler metric, and let
\(\tilde{\omega} = \omega + i \partial \bar{\partial} \varphi\) be the
new metric, where \(\varphi\) is a smooth real-valued function called
the Kähler potential. The condition that \(\tilde{\omega}\) is a metric
requires \(\omega + i \partial \bar{\partial} \varphi > 0\). The
equation \(\text{Ric}(\tilde{\omega}) = \Omega\) can be written as:
\[ \det \left( g_{a\bar{b}} + \frac{\partial^2 \varphi}{\partial z^a \partial \bar{z}^b} \right) = e^F \det(g_{a\bar{b}}) \]
where \(F\) is a smooth function determined by \(\omega\) and the target
Ricci form \(\Omega\), normalized such that:
\[ \int_M e^F \omega^n = \int_M \omega^n \] The goal is to prove the
existence of a unique smooth solution \(\varphi\) to this equation.

\subsection{5.2 The Continuity Method}\label{ch1982shingtungyau:the-continuity-method}

To solve this fully nonlinear PDE, Yau employed the continuity method.
He introduced a parameter \(t \in [0,1]\) and considered the family of
equations:
\[ \det \left( g_{a\bar{b}} + \frac{\partial^2 \varphi_t}{\partial z^a \partial \bar{z}^b} \right) = e^{tF + c_t} \det(g_{a\bar{b}}) \]
where \(c_t\) is a constant chosen to satisfy the volume normalization.
Let \(S \subset [0,1]\) be the set of \(t\) for which a smooth solution
\(\varphi_t\) exists. The goal is to show that \(S\) is non-empty, open,
and closed: - \textbf{Non-emptiness:} At \(t=0\), \(\varphi_0 = 0\) is a
trivial solution, so \(0 \in S\). - \textbf{Openness:} Follows from the
implicit function theorem. The linearized operator of the Monge-Ampère
equation is the Laplace-Beltrami operator associated with the metric
\(g_{t}\). Since \(g_t\) is Kähler, the Laplacian is elliptic, and its
kernel consists of constants. - \textbf{Closedness:} Requires proving a
priori estimates on \(\varphi_t\) and its derivatives up to third order,
independent of \(t\). If these estimates hold, any sequence of solutions
\(\varphi_{t_k}\) as \(t_k \to t\) will converge to a smooth solution
\(\varphi_t\).

\subsection{\texorpdfstring{5.3 The \(C^0\) Estimate via Moser
Iteration}{5.3 The C\^{}0 Estimate via Moser Iteration}}\label{ch1982shingtungyau:the-c0-estimate-via-moser-iteration}

The most difficult step of the proof was establishing the \(C^0\)
estimate: \[ \sup_M |\varphi_t| \le C \] for a constant \(C\)
independent of \(t\). Because the equation is fully nonlinear, classical
elliptic estimates could not be applied.

Yau developed an innovative technique that combined integration by parts
with Moser iteration. By testing the Monge-Ampère equation against
powers of \(\varphi\) and utilizing Sobolev inequalities, Yau showed
that the \(L^p\) norms of \(\varphi\) could be iteratively bounded,
eventually yielding a uniform bound on the supremum norm of \(\varphi\).

\subsection{5.4 Higher-Order Estimates and
Regularity}\label{ch1982shingtungyau:higher-order-estimates-and-regularity}

Once the \(C^0\) estimate was established, Yau proved the \(C^2\) and
\(C^3\) estimates: 1. \textbf{The \(C^2\) Estimate:} Yau bounded the
Laplacian of \(\varphi\) (which controls the second derivatives) by
applying the maximum principle to a carefully chosen auxiliary function
involving the curvature tensor of the background metric. 2. \textbf{The
\(C^3\) Estimate:} Using the \(C^0\) and \(C^2\) estimates, Yau
established bounds on the third derivatives of \(\varphi\) by adapting a
calculation originally due to Calabi. 3. \textbf{Schauder Estimates:}
With the third derivatives bounded, the equation became uniformly
elliptic, allowing the application of Evans-Krylov theory and classical
Schauder estimates to obtain \(C^{k,\alpha}\) bounds for all
\(k \ge 2\). This completed the proof of closedness, establishing the
existence of a unique smooth solution.

\subsection{5.5 Geometric Implications of the
Proof}\label{ch1982shingtungyau:geometric-implications-of-the-proof}

The proof of the complex Monge-Ampère equation had immediate,
far-reaching geometric implications. It provided the first known
non-trivial compact Riemannian manifolds with holonomy group equal to
\(SU(n)\).

Before Yau's proof, the only known manifolds with Ricci-flat metrics
were flat tori. By proving that compact manifolds with zero first Chern
class admit Ricci-flat metrics, Yau opened a new field of study in
Riemannian geometry, showing that complex geometry can dictate the
metric properties of a manifold.

\begin{center}\rule{0.5\linewidth}{0.5pt}\end{center}

\section{6. Calabi-Yau Manifolds}\label{ch1982shingtungyau:calabi-yau-manifolds}

\subsection{6.1 Definition and Holonomy
Group}\label{ch1982shingtungyau:definition-and-holonomy-group}

A compact Kähler manifold \(M\) with vanishing first Chern class
\(c_1(M) = 0\) is called a Calabi-Yau manifold. Yau's proof of the
Calabi conjecture implies that every Calabi-Yau manifold admits a unique
Ricci-flat Kähler metric in each Kähler class. This metric has vanishing
Ricci curvature (\(\text{Ric} = 0\)).

In the language of holonomy groups, a Calabi-Yau manifold of complex
dimension \(n\) has a holonomy group contained within the special
unitary group \(SU(n)\). The presence of \(SU(n)\) holonomy is
equivalent to the existence of a globally defined, non-vanishing,
covariantly constant holomorphic volume form.

\subsection{6.2 Topological and Geometric
Properties}\label{ch1982shingtungyau:topological-and-geometric-properties}

Calabi-Yau manifolds possess several notable properties: - The Hodge
numbers satisfy the symmetry \(h^{p,q} = h^{n-p, q}\) and
\(h^{p,0} = 0\) for \(0 < p < n\), while \(h^{0,0} = h^{n,0} = 1\). -
They admit a unique holomorphic \(n\)-form \(\Omega\), which is
covariantly constant with respect to the Ricci-flat metric. - By Yau's
theorem, their deformation space is unobstructed, as shown by Bogomolov,
Tian, and Todorov. This means that the moduli space of complex
structures on a Calabi-Yau manifold is a smooth manifold of dimension
\(h^{n-1, 1}\).

\subsection{6.3 Key Examples: K3 Surfaces and Quintic
Threefolds}\label{ch1982shingtungyau:key-examples-k3-surfaces-and-quintic-threefolds}

Calabi-Yau manifolds exist in every complex dimension: - \textbf{Complex
Dimension 1:} The only compact Calabi-Yau manifold of dimension 1 is a
complex 1-torus, which is topologically a two-dimensional torus
(elliptic curve). - \textbf{Complex Dimension 2:} The only compact
Calabi-Yau manifolds of dimension 2 are K3 surfaces and complex 2-tori.
K3 surfaces are topologically unique but admit a 20-dimensional family
of hyperkähler metrics. - \textbf{Complex Dimension 3:} The well-known
example is the quintic threefold in complex projective space
\(\mathbb{C}\mathbb{P}^4\). This manifold is defined by the vanishing of
a homogeneous polynomial of degree 5, such as the Fermat quintic:
\[ z_0^5 + z_1^5 + z_2^5 + z_3^5 + z_4^5 = 0 \] Topologically, it has
Euler characteristic \(\chi = -200\), and its Hodge numbers are
\(h^{1,1} = 1\) and \(h^{2,1} = 101\).

\subsection{6.4 Mathematical Applications to Algebraic
Geometry}\label{ch1982shingtungyau:mathematical-applications-to-algebraic-geometry}

In algebraic geometry, Calabi-Yau manifolds represent one of the three
fundamental building blocks in the classification of higher-dimensional
algebraic varieties. According to the Minimal Model Program, any complex
projective variety can be built from: - Fano varieties (which have
positive first Chern class). - Calabi-Yau varieties (which have zero
first Chern class). - Varieties of general type (which have negative
first Chern class).

Yau's theorem proved that Calabi-Yau varieties are exactly those that
admit flat connections on their tangent bundles, establishing a link
between algebraic classification and differential geometry.

\subsection{6.5 The Strominger-Yau-Zaslow (SYZ)
Geometry}\label{ch1982shingtungyau:the-strominger-yau-zaslow-syz-geometry}

In 1996, Strominger, Yau, and Zaslow introduced a geometric framework to
describe the boundary behavior of mirror symmetry. They proposed that
Calabi-Yau manifolds admit fibrations by special Lagrangian tori.

This fibration structure allows one to view the manifold as a torus
bundle over a base space. Mirror symmetry is then realized by replacing
each fiber torus with its dual torus. This formulation has provided
algebraic and symplectic geometers with a concrete geometric picture of
mirror symmetry.

\begin{center}\rule{0.5\linewidth}{0.5pt}\end{center}

\section{7. The Positive Mass Conjecture:
Background}\label{ch1982shingtungyau:the-positive-mass-conjecture-background}

\subsection{7.1 Gravity and Curvature in General
Relativity}\label{ch1982shingtungyau:gravity-and-curvature-in-general-relativity}

In Albert Einstein's general theory of relativity, spacetime is modeled
as a 4-dimensional Lorentzian manifold \((N, g)\) whose curvature is
related to the energy and momentum of matter via the Einstein field
equations:
\[ G_{\mu\nu} = R_{\mu\nu} - \frac{1}{2}Rg_{\mu\nu} = 8\pi T_{\mu\nu} \]
where \(G_{\mu\nu}\) is the Einstein tensor and \(T_{\mu\nu}\) is the
stress-energy tensor.

A fundamental physical principle is that the total energy of an isolated
system must be non-negative. Without this positivity, gravity would be
unstable: a vacuum could spontaneously decay into positive and negative
energy states, leading to runaway collapse.

\subsection{7.2 The Concept of ADM Mass}\label{ch1982shingtungyau:the-concept-of-adm-mass}

Because of the equivalence principle, gravity cannot be localized to a
single point, meaning there is no local mass density. However, for an
isolated gravitational system (like a star or a black hole), the
spacetime becomes flat at large distances.

We model this as an asymptotically flat initial data set \((M, g, K)\),
which represents a spatial slice of spacetime. The total energy and
momentum of this slice can be measured by boundary integrals at spatial
infinity, defined by Richard Arnowitt, Stanley Deser, and Charles Misner
(ADM) in 1959. For a 3-manifold \(M\) with metric \(g_{ij}\) asymptotic
to the flat metric \(\delta_{ij}\), the ADM mass \(E\) is defined as:
\[ E = \frac{1}{16\pi} \lim_{r \to \infty} \int_{S_r} \sum_{i,j=1}^3 \left( \frac{\partial g_{ij}}{\partial x_i} - \frac{\partial g_{ii}}{\partial x_j} \right) \nu_j \, d\sigma \]
where \(S_r\) is a coordinate sphere of radius \(r\) and \(\nu\) is the
outward unit normal.

\subsection{7.3 Statement of the Conjecture and Energy
Conditions}\label{ch1982shingtungyau:statement-of-the-conjecture-and-energy-conditions}

The positive mass conjecture states that if \((M, g, K)\) is an
asymptotically flat initial data set satisfying the Einstein constraint
equations: \[ R(g) - |K|^2_g + (\text{tr}_g K)^2 = 16\pi \mu \]
\[ \text{div}_g (K - (\text{tr}_g K)g) = 8\pi J \] and the dominant
energy condition (which states that the local energy density \(\mu\) is
at least as large as the local momentum density \(J\),
\(\mu \ge |J|_g\)), then the ADM mass \(E\) and the ADM momentum \(P\)
must satisfy: \[ E \ge |P| \] In particular, \(E \ge 0\), and \(E = 0\)
if and only if \(M\) is a slice of flat Minkowski spacetime. Proving
this global inequality from local constraints was one of the major
challenges of mathematical relativity.

\begin{center}\rule{0.5\linewidth}{0.5pt}\end{center}

\section{8. The Positive Mass Conjecture: Minimal Surface
Proof}\label{ch1982shingtungyau:the-positive-mass-conjecture-minimal-surface-proof}

\subsection{8.1 Schoen and Yau's
Collaboration}\label{ch1982shingtungyau:schoen-and-yaus-collaboration}

In 1979, Shing-Tung Yau and his former PhD student Richard Schoen
published a proof of the positive mass conjecture for three-dimensional
manifolds. Their proof used the theory of minimal surfaces---surfaces
that minimize area and have zero mean curvature---in a completely novel
way.

Rather than studying the equations of motion directly, they used minimal
surfaces as topological and geometric probes to explore the global
structure of the manifold.

\subsection{8.2 The Time-Symmetric Case and Conformal
Deformation}\label{ch1982shingtungyau:the-time-symmetric-case-and-conformal-deformation}

The proof is most intuitive in the time-symmetric case, where the
extrinsic curvature \(K\) vanishes. The dominant energy condition then
simplifies to the statement that the scalar curvature \(R\) of the
3-manifold \(M\) is non-negative: \[ R \ge 0 \] Schoen and Yau argued by
contradiction. They assumed that the ADM mass \(E\) was negative. Under
this assumption, they conformal-deformed the metric to \(u^4 g\), where
\(u\) is a positive function asymptotic to 1 at infinity, such that the
deformed metric had zero scalar curvature and negative ADM mass.

This allowed them to construct a stable minimal surface \(\Sigma\) of
codimension 1 (a 2-dimensional surface in the 3-manifold) that extends
to infinity. The stability of \(\Sigma\) means that the second variation
of its area is non-negative:
\[ \int_{\Sigma} \left( |\nabla \psi|^2 + \frac{1}{2}(R - K_{\Sigma} + |A|^2)\psi^2 \right) \ge 0 \]
for all test functions \(\psi\), where \(K_{\Sigma}\) is the Gauss
curvature of \(\Sigma\) and \(|A|^2\) is the norm of its second
fundamental form. By choosing \(\psi\) appropriately and applying the
Gauss-Bonnet theorem, they showed that the Gauss curvature would
integrate to a contradiction, proving that \(E \ge 0\).

\subsection{8.3 The Jang Equation and the General
Case}\label{ch1982shingtungyau:the-jang-equation-and-the-general-case}

For the general case where the extrinsic curvature \(K\) is non-zero,
Schoen and Yau analyzed the Jang equation. The Jang equation is a
quasi-linear elliptic PDE for a function \(f: M \to \mathbb{R}\):
\[ \sum_{i,j=1}^3 \left( g^{ij} - \frac{f^i f^j}{1 + |\nabla f|^2} \right) \left( \nabla_i \nabla_j f - K_{ij} \right) = 0 \]
where \(f^i = g^{ik} \partial_k f\).

The solution \(f\) represents a graph
\(\Sigma_f \subset M \times \mathbb{R}\). The scalar curvature of the
induced metric on \(\Sigma_f\) is positive due to the dominant energy
condition. By establishing the existence and regularity of solutions to
the Jang equation, Schoen and Yau reduced the general case to the
time-symmetric case.

\subsection{8.4 Singularity Analysis and Dimensional
Generalization}\label{ch1982shingtungyau:singularity-analysis-and-dimensional-generalization}

Schoen and Yau's minimal surface proof applied directly to manifolds of
dimension \(n \le 7\). In dimensions \(n \ge 8\), minimal hypersurfaces
can develop singularities, which prevents the direct application of
their regularity results.

Schoen and Yau later developed techniques to resolve these
singularities, extending their minimal surface proof to all dimensions.
In 1981, Edward Witten provided a different proof using the Dirac
operator and spinors, which applies to spin manifolds of any dimension.
However, Schoen and Yau's minimal surface proof remains a cornerstone of
geometric analysis, demonstrating how variational methods can solve
global problems in physical geometry.

\begin{center}\rule{0.5\linewidth}{0.5pt}\end{center}

\section{9. Real Monge-Ampere
Equations}\label{ch1982shingtungyau:real-monge-ampere-equations}

\subsection{9.1 Formulation and
Ellipticity}\label{ch1982shingtungyau:formulation-and-ellipticity}

The real Monge-Ampère equation is a fully nonlinear second-order
elliptic partial differential equation. For a function \(u\) defined on
a domain \(\Omega \subset \mathbb{R}^n\), it is formulated as:
\[ \det(D^2 u) = f(x, u, \nabla u) \] where
\(D^2 u = \left( \frac{\partial^2 u}{\partial x_i \partial x_j} \right)\)
is the Hessian matrix of \(u\). The equation is elliptic only when \(u\)
is convex, which ensures that all eigenvalues of \(D^2 u\) are positive.
The real Monge-Ampère equation arises naturally in optimal transport,
differential geometry (such as the Minkowski problem), and meteorology.

\subsection{9.2 Cheng and Yau's Existence
Theorems}\label{ch1982shingtungyau:cheng-and-yaus-existence-theorems}

Shing-Tung Yau and Shiu-Yuen Cheng studied the Dirichlet problem for the
real Monge-Ampère equation on convex domains \(\Omega\):
\[ \det(D^2 u) = f(x) \quad \text{in } \Omega, \quad u = g \quad \text{on } \partial \Omega \]
They proved the existence, uniqueness, and regularity of classical
solutions when the boundary \(\partial \Omega\) is strictly convex and
\(f(x) > 0\) is smooth. Their proof required establishing uniform bounds
on the Hessian and higher-order derivatives of \(u\) using the maximum
principle and barrier functions.

\subsection{9.3 Applications to Affine Differential
Geometry}\label{ch1982shingtungyau:applications-to-affine-differential-geometry}

Cheng and Yau used their analysis of the real Monge-Ampère equation to
solve major problems in affine differential geometry, which studies the
properties of manifolds that are invariant under affine transformations.

They resolved the affine Bernstein conjecture, which was proposed by
Shiing-Shen Chern. The conjecture asks whether a complete affine maximal
hypersurface in \(\mathbb{R}^{n+1}\) must be a paraboloid. By
constructing complete affine Kähler metrics using solutions to the
Monge-Ampère equation, Cheng and Yau proved the conjecture for all
dimensions, establishing a deep connection between convex analysis and
affine geometry.

\begin{center}\rule{0.5\linewidth}{0.5pt}\end{center}

\section{10. Complex Monge-Ampere
Equations}\label{ch1982shingtungyau:complex-monge-ampere-equations}

\subsection{10.1 Non-compact Kähler
Manifolds}\label{ch1982shingtungyau:non-compact-kuxe4hler-manifolds}

In addition to his work on compact Kähler manifolds, Yau studied the
complex Monge-Ampère equation on non-compact manifolds. On a complex
manifold \(M\) of dimension \(n\), the complex Monge-Ampère equation is
written as:
\[ (\omega + i \partial \bar{\partial} \varphi)^n = e^F \omega^n \] If
\(M\) is non-compact, the behavior of the solution \(\varphi\) at
infinity must be controlled. The study of non-compact complex
Monge-Ampère equations is crucial for constructing complete
Kähler-Einstein metrics on open complex manifolds.

\subsection{10.2 Complete Kähler-Einstein Metrics on Bounded
Domains}\label{ch1982shingtungyau:complete-kuxe4hler-einstein-metrics-on-bounded-domains}

In collaboration with Shiu-Yuen Cheng, Yau investigated the existence of
complete Kähler-Einstein metrics on bounded domains
\(\Omega \subset \mathbb{C}^n\). A Kähler-Einstein metric \(g\)
satisfies \(\text{Ric}(g) = \lambda g\). For \(\lambda < 0\), the
equation for the Kähler potential \(\varphi\) becomes a complex
Monge-Ampère equation that blows up at the boundary:
\[ \det \left( \frac{\partial^2 \varphi}{\partial z_i \partial \bar{z}_j} \right) = e^{(n+1)\varphi} \]
Cheng and Yau proved that if \(\Omega\) is a strictly pseudoconvex
domain, there exists a unique complete Kähler-Einstein metric on
\(\Omega\) with negative scalar curvature.

\subsection{10.3 Dirichlet Problem at
Infinity}\label{ch1982shingtungyau:dirichlet-problem-at-infinity}

Cheng and Yau also resolved the Dirichlet problem at infinity for the
complex Monge-Ampère equation. They proved that one can prescribe the
boundary behavior of the Kähler potential at the boundary of a
pseudoconvex domain.

This work generalized Charles Fefferman's asymptotic expansions of the
Bergman metric and provided a complete geometric description of the
boundary behavior of Kähler-Einstein metrics, opening new directions in
multi-variable complex analysis.

\begin{center}\rule{0.5\linewidth}{0.5pt}\end{center}

\section{11. The Yamabe Problem and Kahler-Einstein
Metrics}\label{ch1982shingtungyau:the-yamabe-problem-and-kahler-einstein-metrics}

\subsection{11.1 The Yamabe Conjecture and Critical Sobolev
Exponents}\label{ch1982shingtungyau:the-yamabe-conjecture-and-critical-sobolev-exponents}

The Yamabe problem asks whether every compact Riemannian manifold
\((M, g)\) of dimension \(n \ge 3\) can be conformally deformed to a
metric with constant scalar curvature. This is equivalent to finding a
positive function \(u\) satisfying the semi-linear elliptic PDE:
\[ -\frac{4(n-1)}{n-2} \Delta_g u + R_g u = \tilde{R} u^{\frac{n+2}{n-2}} \]
where \(\tilde{R}\) is a constant. The exponent \(\frac{n+2}{n-2}\) is
the critical Sobolev exponent, which makes the equation difficult to
solve because the associated variational functional lacks compactness.

\subsection{11.2 Schoen's Resolution and Positive Mass
Theorem}\label{ch1982shingtungyau:schoens-resolution-and-positive-mass-theorem}

The Yamabe conjecture was resolved by Richard Schoen in 1984. His proof
relied on the positive mass theorem that he had developed with Yau.

Schoen showed that the conformal invariant (the Yamabe constant) of any
compact manifold not conformally equivalent to the sphere is strictly
less than that of the sphere. He achieved this by analyzing the Green's
function of the conformal Laplacian and showing that its expansion at
the pole is controlled by the ADM mass. Because the mass is positive,
the expansion yields a strict inequality, establishing the existence of
a minimizer for the Yamabe functional.

\subsection{11.3 Kähler-Einstein Metrics with Signed Chern
Classes}\label{ch1982shingtungyau:kuxe4hler-einstein-metrics-with-signed-chern-classes}

The existence of Kähler-Einstein metrics (\(\text{Ric}(g) = \lambda g\))
on compact Kähler manifolds depends on the sign of the first Chern class
\(c_1(M)\): - \textbf{Negative First Chern Class (\(c_1(M) < 0\)):}
Thierry Aubin and Shing-Tung Yau independently proved in 1976 that there
always exists a unique Kähler-Einstein metric with negative scalar
curvature (\(\lambda < 0\)). - \textbf{Zero First Chern Class
(\(c_1(M) = 0\)):} Yau's proof of the Calabi conjecture established the
existence of a unique Ricci-flat (\(\lambda = 0\)) Kähler-Einstein
metric. - \textbf{Positive First Chern Class (\(c_1(M) > 0\)):} Known as
Fano manifolds. The existence of Kähler-Einstein metrics is obstructed
by the Futaki invariant. The resolution of this case was completed in
2015 by Xiuxiong Chen, Simon Donaldson, and Song Sun, who proved that
the existence of a metric is equivalent to K-stability.

\begin{center}\rule{0.5\linewidth}{0.5pt}\end{center}

\section{12. The Smith Conjecture}\label{ch1982shingtungyau:the-smith-conjecture}

\subsection{12.1 Statement of the Smith
Conjecture}\label{ch1982shingtungyau:statement-of-the-smith-conjecture}

The Smith conjecture is a classic topological problem formulated by Paul
A. Smith. It states that if a cyclic group \(\mathbb{Z}_p\) of prime
order \(p\) acts differentiably on the 3-sphere \(S^3\) such that the
fixed point set of the action is a 1-dimensional manifold (which must be
homeomorphic to a circle), then this fixed point set cannot be a knotted
circle. In other words, any differentiable periodic action on \(S^3\)
with a circle of fixed points must be conjugate to a standard rotation.

\subsection{12.2 Equivariant Loop Theorem and Dehn's
Lemma}\label{ch1982shingtungyau:equivariant-loop-theorem-and-dehns-lemma}

In 1980, a group of mathematicians, including Shing-Tung Yau, William
Meeks III, William Thurston, Hyman Bass, and Peter Shalen, proved the
Smith conjecture.

Meeks and Yau's contribution was the application of minimal surface
theory to establish the equivariant loop theorem and Dehn's lemma. They
showed that if a 3-manifold admits a group action, one can construct
minimal disks that are invariant under the group. This allowed them to
perform topological surgery while preserving the symmetry of the group
action.

\subsection{12.3 Minimal Surfaces and Thurston's
Geometrization}\label{ch1982shingtungyau:minimal-surfaces-and-thurstons-geometrization}

By using Meeks and Yau's equivariant Dehn's lemma, the topological
problem of the Smith conjecture was reduced to showing that the
universal cover of a Haken manifold is geometric. This was resolved
using William Thurston's geometrization theorem. The proof of the Smith
conjecture was a major triumph, demonstrating that the analytical
properties of minimal surfaces could resolve long-standing topological
conjectures.

\begin{center}\rule{0.5\linewidth}{0.5pt}\end{center}

\section{13. Differential Geometry of
Submanifolds}\label{ch1982shingtungyau:differential-geometry-of-submanifolds}

\subsection{13.1 Minimal Submanifolds and
Stability}\label{ch1982shingtungyau:minimal-submanifolds-and-stability}

Minimal submanifolds are submanifolds whose mean curvature vanishes,
meaning they are critical points of the volume functional. A minimal
submanifold is stable if the second variation of its volume is
non-negative for all variations.

Shing-Tung Yau studied the stability of minimal submanifolds and their
relationship to the topology of the ambient manifold. In collaboration
with William Meeks, Yau proved that the boundary of a stable minimal
surface in a 3-manifold satisfies strict topological constraints, which
has applications in three-dimensional topology.

\subsection{13.2 The Li-Yau Gradient
Estimates}\label{ch1982shingtungyau:the-li-yau-gradient-estimates}

In 1986, Shing-Tung Yau and Peter Li published a paper containing what
is now known as the Li-Yau gradient estimate for the heat equation. For
a positive solution \(u(x,t)\) of the heat equation:
\[ \left( \frac{\partial}{\partial t} - \Delta \right) u = 0 \] on a
complete Riemannian manifold with Ricci curvature bounded from below by
\(-K\), they proved:
\[ \frac{|\nabla u|^2}{u^2} - \alpha \frac{u_t}{u} \le \frac{n \alpha^2}{2t} + \frac{n \alpha^2 K}{2(\alpha - 1)} \]
for any \(\alpha > 1\). This estimate implies that the gradient of the
log of the solution is controlled by the value of the solution itself
and the geometry of the manifold. The Li-Yau estimate is a fundamental
tool for obtaining heat kernel bounds, Harnack inequalities, and
eigenvalue estimates.

\subsection{13.3 Yau's Conjecture on
Eigenvalues}\label{ch1982shingtungyau:yaus-conjecture-on-eigenvalues}

Yau formulated several conjectures regarding the eigenvalues of the
Laplace-Beltrami operator on manifolds. One of the well-known is Yau's
conjecture on the first eigenvalue of minimal hypersurfaces in spheres.

The conjecture states that if \(\Sigma^n\) is a closed minimal
hypersurface in the unit sphere \(S^{n+1}(1)\), then the first
eigenvalue \(\lambda_1(\Sigma)\) of the Laplacian on \(\Sigma\) is
exactly \(n\). While this conjecture has been proven for specific cases
(such as when the hypersurface is symmetric), it remains open in its
general formulation.

\begin{center}\rule{0.5\linewidth}{0.5pt}\end{center}

\section{14. Black Hole Formation and General
Relativity}\label{ch1982shingtungyau:black-hole-formation-and-general-relativity}

\subsection{14.1 Apparent Horizons and Gravity
Collapse}\label{ch1982shingtungyau:apparent-horizons-and-gravity-collapse}

In general relativity, a black hole is a region of spacetime from which
nothing can escape. The boundary of a black hole is the event horizon,
which is defined globally. Because event horizons are difficult to
locate on a given spatial slice, physicists use the local notion of an
apparent horizon. An apparent horizon is a trapped surface where the
expansion of outgoing light rays vanishes, meaning gravity is strong
enough to bend light inward.

\subsection{14.2 The Schoen-Yau Limit on Matter
Concentration}\label{ch1982shingtungyau:the-schoen-yau-limit-on-matter-concentration}

In 1983, Richard Schoen and Shing-Tung Yau proved a mathematical theorem
regarding black hole formation. They showed that if a given region of
space contains a sufficiently high concentration of matter, then the
gravitational force must be so strong that a trapped surface (and thus a
black hole) must form.

Specifically, they proved that if the matter density \(\rho\) and the
size of the region \(U\) (measured by the radius \(R\) of the largest
inscribed torus) satisfy: \[ R^2 \rho \ge C \] where \(C\) is a
universal constant, then \(U\) must contain a trapped surface. This was
the first rigorous mathematical proof that matter concentration leads to
gravitational collapse, confirming a central prediction of general
relativity.

\subsection{14.3 The Wang-Yau Quasi-local
Mass}\label{ch1982shingtungyau:the-wang-yau-quasi-local-mass}

In general relativity, there is no local mass density, which makes it
difficult to define the mass contained in a finite region of space. In
2009, Shing-Tung Yau and Mu-Tao Wang defined the Wang-Yau quasi-local
mass.

Their definition is based on embedding the boundary of a region (a
topological 2-sphere) into reference spacetimes (such as Minkowski
space) and minimizing a Hamiltonian energy functional. The Wang-Yau
quasi-local mass satisfies all the required physical properties,
including positivity, and has been used to study gravitational waves and
black hole dynamics.

\subsection{14.4 Apparent Horizon vs Event
Horizon}\label{ch1982shingtungyau:apparent-horizon-vs-event-horizon}

Schoen and Yau's concentration theorems focused on apparent horizons
because they are determined by local geometry. This is computationally
advantageous, as event horizons require knowing the complete future
development of the spacetime.

Their work provided the mathematical foundation for numerical
relativists simulating the collision of black holes. It proved that if
the simulation code detects a region satisfying their density
inequality, a black hole boundary is guaranteed to exist.

\begin{center}\rule{0.5\linewidth}{0.5pt}\end{center}

\section{15. Mirror Symmetry and String
Theory}\label{ch1982shingtungyau:mirror-symmetry-and-string-theory}

\subsection{15.1 Compactification and the Extra
Dimensions}\label{ch1982shingtungyau:compactification-and-the-extra-dimensions}

In superstring theory, the fundamental building blocks of matter are
1-dimensional strings vibrating in a 10-dimensional spacetime. To
reconcile this with our 4-dimensional experience, the extra 6 dimensions
must be compactified into an extremely small space.

For the 4-dimensional physics to preserve \(\mathcal{N}=1\)
supersymmetry, this 6-dimensional space must be a complex 3-dimensional
Calabi-Yau manifold. The shape and topology of the Calabi-Yau manifold
determine the properties of the particles (such as their masses and
charges) in the 4-dimensional universe.

\subsection{15.2 Mirror Symmetry as a Physical
Duality}\label{ch1982shingtungyau:mirror-symmetry-as-a-physical-duality}

Physicists discovered that two different string theories (Type IIA and
Type IIB) compactified on two geometrically distinct Calabi-Yau
manifolds \(X\) and \(Y\) can result in the same 4-dimensional physical
theory. In this case, \(X\) and \(Y\) are called a mirror pair.

Mirror symmetry relates the symplectic geometry of \(X\) (specifically,
the A-model topological string theory, which involves counting
holomorphic curves) to the complex geometry of \(Y\) (the B-model
theory, which involves classical period integrals). In 1991, physicists
used mirror symmetry to calculate the number of rational curves of any
degree on a quintic threefold, solving a problem in algebraic geometry
that had resisted mathematicians for over a century.

\subsection{15.3 The SYZ Conjecture: Special Lagrangian
Fibrations}\label{ch1982shingtungyau:the-syz-conjecture-special-lagrangian-fibrations}

To provide a geometric explanation for mirror symmetry, Andrew
Strominger, Shing-Tung Yau, and Eric Zaslow proposed the SYZ conjecture
in 1996. They conjectured that a Calabi-Yau threefold \(X\) should admit
a fibration by special Lagrangian tori: \[ \pi: X \to B \] and that its
mirror manifold \(Y\) is obtained by taking the dual torus fibration
over the same base \(B\). Mirror symmetry is then understood as
T-duality along these torus fibers. The SYZ conjecture remains a primary
guiding principle for mathematicians working to understand the geometric
mechanics of mirror symmetry.

\subsection{15.4 Gromov-Witten Invariants and Curve
Counting}\label{ch1982shingtungyau:gromov-witten-invariants-and-curve-counting}

Mirror symmetry allowed mathematicians to compute Gromov-Witten
invariants, which count the number of holomorphic curves on a manifold.
On a quintic threefold, this involves finding the number of rational
curves of degree \(d\).

For degree 1 (lines), the number is 2875, a result known since the
nineteenth century. For degree 2 (conics), the number is 609250. Prior
to mirror symmetry, the degree 3 case was unsolved. Mirror symmetry
provided a generating function that computed these numbers for all
degrees simultaneously, a result that was later proven rigorously using
localization techniques in algebraic geometry.

\begin{center}\rule{0.5\linewidth}{0.5pt}\end{center}

\section{16. Unifying Techniques in Geometric
Analysis}\label{ch1982shingtungyau:unifying-techniques-in-geometric-analysis}

\subsection{16.1 Yau's Maximum Principle on Non-compact
Manifolds}\label{ch1982shingtungyau:yaus-maximum-principle-on-non-compact-manifolds}

One of Shing-Tung Yau's effective analytical contributions is his
maximum principle for complete, non-compact Riemannian manifolds. The
classical maximum principle states that a subharmonic function on a
compact domain must achieve its maximum on the boundary.

On a non-compact manifold, a function may not achieve its maximum. Yau
proved that if \(u\) is a smooth function bounded from above on a
complete Riemannian manifold whose Ricci curvature is bounded from
below, then there exists a sequence of points \(x_k\) such that:
\[ \lim_{k\to\infty} u(x_k) = \sup u, \quad \lim_{k\to\infty} |\nabla u(x_k)| = 0, \quad \lim_{k\to\infty} \Delta u(x_k) \le 0 \]
This result allows mathematicians to apply maximum principle arguments
to open manifolds, which is crucial for proving uniqueness and existence
theorems for geometric PDEs.

\subsection{16.2 Differential Harnack
Inequalities}\label{ch1982shingtungyau:differential-harnack-inequalities}

Yau and Peter Li developed differential Harnack inequalities for the
heat equation. By bounding the gradient of a solution in terms of the
solution itself, they showed that one can compare the values of the
solution at different points in space and time.

These estimates were later adapted by Richard Hamilton to the Ricci
flow, a parabolic PDE that deforms a Riemannian metric in the direction
of its Ricci curvature:
\[ \frac{\partial}{\partial t} g_{ij} = -2 R_{ij} \] Hamilton's
differential Harnack estimate for the Ricci flow was a crucial step in
the analysis of singularities.

\subsection{16.3 Connections to the Ricci Flow and Poincaré
Conjecture}\label{ch1982shingtungyau:connections-to-the-ricci-flow-and-poincaruxe9-conjecture}

The analytical tools developed by Yau laid the foundation for the Ricci
flow program. Grigori Perelman's proof of the Poincaré conjecture and
Thurston's geometrization conjecture relied on the analysis of
singularities in the Ricci flow, which built directly on the gradient
estimates, maximum principles, and blow-up analysis developed by Yau and
his collaborators. Yau's work demonstrated that nonlinear diffusion
equations can be used to decompose and classify topological spaces.

\begin{center}\rule{0.5\linewidth}{0.5pt}\end{center}

\section{17. Broad Impact on Modern
Mathematics}\label{ch1982shingtungyau:broad-impact-on-modern-mathematics}

\subsection{17.1 Mentorship and PhD
Students}\label{ch1982shingtungyau:mentorship-and-phd-students}

Shing-Tung Yau has had a significant impact on modern mathematics through
his mentorship. Over his career, he has supervised more than 70 PhD
students, many of whom have become leaders in geometric analysis,
differential geometry, and mathematical physics.

His students include Richard Schoen, Shiu-Yuen Cheng, Peter Li, Mu-Tao
Wang, and Kefeng Liu. Through his active seminars and collaborations,
Yau has created a global network of researchers dedicated to the study
of geometric analysis.

\subsection{17.2 The 120 Open Problems in
Geometry}\label{ch1982shingtungyau:the-120-open-problems-in-geometry}

In 1982, Yau organized a seminar on differential geometry at the
Institute for Advanced Study, which resulted in the publication of the
book \emph{Seminar on Differential Geometry}. In the introduction, Yau
compiled a list of 120 open problems in geometry and analysis.

This list served as a roadmap for the development of geometric analysis
for the next forty years. Many of the problems led to major mathematical
discoveries, and several became fields of study in their own right,
defining research programs for decades.

\subsection{17.3 Fostering Mathematics in Greater
China}\label{ch1982shingtungyau:fostering-mathematics-in-greater-china}

Yau has been a key builder of mathematical institutions in the
Chinese-speaking world. He helped establish several world-class research
centers: - \textbf{The Institute of Mathematical Sciences} at the
Chinese University of Hong Kong. - \textbf{The Morningside Center of
Mathematics} in Beijing. - \textbf{The Center of Mathematical Sciences}
at Zhejiang University in Hangzhou. - \textbf{The Yau Mathematical
Sciences Center (YMSC)} at Tsinghua University in Beijing.

He also founded the Morningside Medal of Mathematics, awarded to
outstanding Chinese mathematicians under the age of 45, and created
various mathematics competitions for middle and high school students to
foster young talent.

\begin{center}\rule{0.5\linewidth}{0.5pt}\end{center}

\section{18. Awards and Professional
Recognition}\label{ch1982shingtungyau:awards-and-professional-recognition}

\subsection{18.1 Major International
Prizes}\label{ch1982shingtungyau:major-international-prizes}

Shing-Tung Yau's contributions have been recognized with numerous
prestigious awards: - \textbf{1982: Fields Medal.} Awarded for his
contributions to differential equations, the Calabi conjecture, the
positive mass conjecture, and the Monge-Ampère equations. -
\textbf{1984: MacArthur Fellowship.} Awarded by the John D. and
Catherine T. MacArthur Foundation. - \textbf{1994: Crafoord Prize.}
Awarded by the Royal Swedish Academy of Sciences for his development of
geometric analysis and his proof of the Calabi conjecture. -
\textbf{1997: United States National Medal of Science.} Presented by
President Bill Clinton for his outstanding contributions to mathematics
and its applications to physics. - \textbf{2010: Wolf Prize in
Mathematics.} Awarded by the Wolf Foundation in Israel ``for his work in
geometric analysis that has had a significant and dramatic impact on many
areas of geometry and physics.'' - \textbf{2018: Marcel Grossmann
Award.} Awarded for his contributions to the mathematical foundations of
general relativity, particularly the positive mass theorem and
quasi-local mass.

\subsection{18.2 Academician and Society
Memberships}\label{ch1982shingtungyau:academician-and-society-memberships}

Yau is an elected member of several prestigious scientific academies,
including: - \textbf{The United States National Academy of Sciences}
(elected 1993). - \textbf{The American Academy of Arts and Sciences}
(elected 1983). - \textbf{The Chinese Academy of Sciences} (elected as a
foreign member in 1994). - \textbf{The Academia Sinica} (elected 1984).

\begin{center}\rule{0.5\linewidth}{0.5pt}\end{center}

\section{19. Selected Major
Publications}\label{ch1982shingtungyau:selected-major-publications}

\subsection{19.1 Groundbreaking Papers in Analysis and
Geometry}\label{ch1982shingtungyau:groundbreaking-papers-in-analysis-and-geometry}

Yau has published over 500 papers and books. The following publications
are considered foundational works in geometric analysis: 1. \textbf{Yau,
S. T. (1978). ``On the Ricci curvature of a compact Kähler manifold and
the complex Monge-Ampère equation, I.'' \emph{Communications on Pure and
Applied Mathematics}, 31(3), 339--411.} - \emph{Description:} This paper
contains the complete proof of the Calabi conjecture and the resolution
of the associated complex Monge-Ampère equation, establishing the
existence of Ricci-flat Kähler metrics. 2. \textbf{Schoen, R., \& Yau,
S. T. (1979). ``On the proof of the positive mass conjecture in general
relativity.'' \emph{Communications in Mathematical Physics}, 65(1),
45--76.} - \emph{Description:} This paper contains the minimal surface
proof of the positive mass theorem for three-dimensional manifolds under
the time-symmetric assumption. 3. \textbf{Schoen, R., \& Yau, S. T.
(1981). ``Proof of the positive mass theorem. II.'' \emph{Communications
in Mathematical Physics}, 79(2), 231--260.} - \emph{Description:} This
paper extends the positive mass proof to the general case (non-zero
extrinsic curvature) by solving the Jang equation.

\subsection{19.2 Relativistic and String Theoretic
Contributions}\label{ch1982shingtungyau:relativistic-and-string-theoretic-contributions}

\begin{enumerate}
\def\labelenumi{\arabic{enumi}.}
\setcounter{enumi}{3}
\tightlist
\item
 \textbf{Cheng, S. Y., \& Yau, S. T. (1980). ``On the existence of a
 complete Kähler metric on noncompact complex manifolds and the
 regularity of Monge-Ampère's equation.'' \emph{Communications on Pure
 and Applied Mathematics}, 33(4), 507--544.}

 \begin{itemize}
 \tightlist
 \item
 \emph{Description:} This paper establishes the existence of complete
 Kähler-Einstein metrics on bounded strictly pseudoconvex domains.
 \end{itemize}
\item
 \textbf{Meeks, W. H., \& Yau, S. T. (1980). ``Topology of
 three-dimensional manifolds and the embedding problems in minimal
 surface theory.'' \emph{Annals of Mathematics}, 112(3), 441--484.}

 \begin{itemize}
 \tightlist
 \item
 \emph{Description:} This paper applies minimal surface theory to
 solve topological embedding problems, laying the groundwork for the
 proof of the Smith conjecture.
 \end{itemize}
\item
 \textbf{Strominger, A., Yau, S. T., \& Zaslow, E. (1996). ``Mirror
 symmetry is T-duality.'' \emph{Nuclear Physics B}, 479(1-2),
 243--259.}

 \begin{itemize}
 \tightlist
 \item
 \emph{Description:} This paper proposes the SYZ conjecture, offering
 a geometric explanation of mirror symmetry via dual torus
 fibrations.
 \end{itemize}
\item
 \textbf{Wang, M. T., \& Yau, S. T. (2009). ``Isometric embeddings into
 the Minkowski space and new quasi-local mass.'' \emph{Communications
 in Mathematical Physics}, 289(3), 919--942.}

 \begin{itemize}
 \tightlist
 \item
 \emph{Description:} This paper presents the construction of the
 Wang-Yau quasi-local mass, providing a mathematically rigorous and
 physically consistent definition of energy in bounded regions.
 \end{itemize}
\end{enumerate}

\begin{center}\rule{0.5\linewidth}{0.5pt}\end{center}

\section{20. Ongoing Legacy and Future
Directions}\label{ch1982shingtungyau:ongoing-legacy-and-future-directions}

\subsection{20.1 Current Research in Tsinghua and
Beyond}\label{ch1982shingtungyau:current-research-in-tsinghua-and-beyond}

Following his retirement from Harvard, Yau remains active as the
Director of the Yau Mathematical Sciences Center at Tsinghua University.
His current research focuses on extending the boundaries of geometric
analysis, general relativity, and the mathematical foundations of string
theory.

He continues to work on the SYZ conjecture, developing the theory of
special Lagrangian submanifolds to describe mirror symmetry in
non-perturbative regimes. He also works on refining the Wang-Yau
quasi-local mass, aiming to define quasi-local angular momentum and
center of mass in general relativity.

\subsection{20.2 Computational Conformal Geometry and Machine
Learning}\label{ch1982shingtungyau:computational-conformal-geometry-and-machine-learning}

In recent years, Yau has expanded his research to include computational
conformal geometry, which applies conformal mapping techniques to
computer graphics, medical imaging, and computer vision. He has also
explored the use of machine learning to identify and classify Calabi-Yau
manifolds.

By using deep neural networks to compute topological invariants and
search for mirror pairs, Yau is integrating modern machine learning with
classical algebraic geometry. This interdisciplinary approach aims to
automate the search for candidate manifolds in string compactification,
demonstrating his ongoing commitment to expanding the boundaries of
mathematics.

\subsection{20.3 Institutional Leadership and Talent
Development}\label{ch1982shingtungyau:institutional-leadership-and-talent-development}

Through Tsinghua University's Qiuzhen College, which Yau founded, he has
established a specialized curriculum for training elite undergraduate
students in mathematics. He oversees recruitment of international
faculty and organizes summer schools to expose young scholars to active
research fields.

His efforts extend to high school mathematics competitions across China
and other Asian nations. By establishing these platforms, Yau has
created a sustainable pipeline of mathematical talent, ensuring that the
legacy of geometric analysis continues to grow.

\begin{center}\rule{0.5\linewidth}{0.5pt}\end{center}

\section{21. Conclusion}\label{ch1982shingtungyau:conclusion}

\subsection{21.1 Summary of
Contributions}\label{ch1982shingtungyau:summary-of-contributions}

Shing-Tung Yau's mathematical achievements have substantially reshaped
modern geometry and analysis. By proving the Calabi conjecture, he
established the existence of Ricci-flat Kähler metrics, providing the
mathematical playground for superstring theory.

His proof of the positive mass theorem settled a foundational question
regarding the stability of general relativity, and his work on real and
complex Monge-Ampère equations established the analytical tools for
geometric analysis. Yau's work demonstrated that the shape of space and
the equations describing physical fields are closely connected.

\subsection{21.2 The Legacy of Shing-Tung
Yau}\label{ch1982shingtungyau:the-legacy-of-shing-tung-yau}

From his childhood in Hong Kong to his position as a leader in the
global mathematical community, Yau has maintained a deep dedication to
the pursuit of truth. His legacy is reflected in the theorems that bear
his name, the institutions he established, and the generations of
mathematicians he mentored.

Yau's research continues to inspire new discoveries at the intersection
of geometry, analysis, and theoretical physics, ensuring that his
influence will be felt for generations to come.

\subsection{21.3 Geometry and the Philosophy of
Beauty}\label{ch1982shingtungyau:geometry-and-the-philosophy-of-beauty}

Throughout his writings and lectures, Yau has emphasized that geometry
is not merely an abstract set of axioms, but a pursuit of natural
beauty. He often draws comparisons between the structure of a
mathematical proof and classical Chinese poetry, noting that both
require balance, economy of expression, and a deep resonance with
nature.

This philosophical perspective has guided his geometric intuition,
leading him to seek solutions that are not only logically sound but also
geometrically elegant. For Yau, the equations of mathematical physics
are the language through which the universe expresses its intrinsic
symmetry and harmony.

\begin{center}\rule{0.5\linewidth}{0.5pt}\end{center}

\begin{center}\rule{0.5\linewidth}{0.5pt}\end{center}

\section*{Exercises}
\addcontentsline{toc}{section}{Exercises}

\begin{exercise} Verify that $\\mathbb{CP}^n$ admits a K\\'{a}hler-Einstein metric by computing the Ricci curvature. \end{exercise}

\begin{exercise} Show that the Calabi-Yau metric on a K3 surface is Ricci-flat by checking the first Chern class vanishes. \end{exercise}

\begin{exercise} Compute the scalar curvature of the standard sphere $S^n$ with the round metric. \end{exercise}

\begin{exercise} Prove that a complete noncompact manifold with positive scalar curvature has a minimum at infinity. \end{exercise}

\begin{center}\rule{0.5\linewidth}{0.5pt}\end{center}

\section*{Summary}
\addcontentsline{toc}{section}{Summary}

Yau proved the Calabi conjecture, showing that every K\'{a}hler manifold with vanishing first Chern class admits a Ricci-flat metric. This result, with profound applications in string theory and mirror symmetry, together with the positive mass theorem, established the deep connections between differential geometry, PDEs, and mathematical physics.

\section{22. References and Further
Reading}\label{ch1982shingtungyau:references-and-further-reading-1}

\subsection{22.1 Academic Papers}\label{ch1982shingtungyau:academic-papers}

\begin{itemize}
\tightlist
\item
 Yau, S. T. (1978). ``On the Ricci curvature of a compact Kähler
 manifold and the complex Monge-Ampère equation, I''.
 \emph{Communications on Pure and Applied Mathematics}, 31(3),
 339--411.
\item
 Schoen, R., \& Yau, S. T. (1979). ``On the proof of the positive mass
 conjecture in general relativity''. \emph{Communications in
 Mathematical Physics}, 65(1), 45--76.
\item
 Schoen, R., \& Yau, S. T. (1981). ``Proof of the positive mass
 theorem. II''. \emph{Communications in Mathematical Physics}, 79(2),
 231--260.
\item
 Cheng, S. Y., \& Yau, S. T. (1980). ``On the existence of a complete
 Kähler metric on noncompact complex manifolds and the regularity of
 Monge-Ampère's equation''. \emph{Communications on Pure and Applied
 Mathematics}, 33(4), 507--544.
\item
 Meeks, W. H., \& Yau, S. T. (1980). ``Topology of three-dimensional
 manifolds and the embedding problems in minimal surface theory''.
 \emph{Annals of Mathematics}, 112(3), 441--484.
\item
 Strominger, A., Yau, S. T., \& Zaslow, E. (1996). ``Mirror symmetry is
 T-duality''. \emph{Nuclear Physics B}, 479(1-2), 243--259.
\item
 Li, P., \& Yau, S. T. (1986). ``On the parabolic kernel of the
 Schrödinger operator''. \emph{Acta Mathematica}, 156(1), 153--201.
\item
 Wang, M. T., \& Yau, S. T. (2009). ``Isometric embeddings into the
 Minkowski space and new quasi-local mass''. \emph{Communications in
 Mathematical Physics}, 289(3), 919--942.
\item
 Schoen, R., \& Yau, S. T. (1983). ``The existence of a black hole due
 to condensation of matter''. \emph{Communications in Mathematical
 Physics}, 90(4), 575--600.
\item
 Yau, S. T. (1975). ``Harmonic functions on complete Riemannian
 manifolds''. \emph{Communications on Pure and Applied Mathematics},
 28(2), 201--228.
\end{itemize}

\subsection{22.2 Books and Monographs}\label{ch1982shingtungyau:books-and-monographs}

\begin{itemize}
\tightlist
\item
 Yau, S. T., \& Nadis, S. (2010). \emph{The Shape of Inner Space:
 String Theory and the Geometry of the Universe's Hidden Dimensions}.
 Basic Books.
\item
 Yau, S. T., \& Nadis, S. (2019). \emph{The Shape of a Life: One
 Mathematician's Search for the Universe's Hidden Geometry}. Yale
 University Press.
\item
 Cao, H. D., \& Yau, S. T. (Eds.). (1992). \emph{Seminar on
 Differential Geometry}. Annals of Mathematics Studies, Princeton
 University Press.
\item
 Yau, S. T. (Ed.). (2013). \emph{Selected Expository Works of
 Shing-Tung Yau with Commentary}. International Press.
\item
 Gu, C. H., Yau, S. T., \& Hou, J. X. (Eds.). (1998).
 \emph{Differential Geometry and Mathematical Physics}. International
 Press.
\end{itemize}

\subsection{22.3 Surveys and Review
Articles}\label{ch1982shingtungyau:surveys-and-review-articles}

\begin{itemize}
\tightlist
\item
 Bourguignon, J. P. (1981). ``The work of Shing-Tung Yau.''
 \emph{Séminaire Bourbaki}, 33ème année, Vol. 1980/81.

 \begin{itemize}
 \tightlist
 \item
 \emph{Summary:} A comprehensive expository review of Yau's early
 work, focusing on the resolution of the Calabi conjecture and
 complex Monge-Ampère equations.
 \end{itemize}
\item
 Schoen, R. (1983). ``The work of Shing-Tung Yau.'' \emph{Proceedings
 of the International Congress of Mathematicians}, Warsaw.

 \begin{itemize}
 \tightlist
 \item
 \emph{Summary:} Richard Schoen's fields medal laudatio for
 Shing-Tung Yau, highlighting his major breakthroughs in geometric
 analysis and general relativity.
 \end{itemize}
\item
 Uhlenbeck, K. (1988). ``Shing-Tung Yau: The founder of modern
 geometric analysis.'' \emph{American Mathematical Society
 Translations}, 141, 12--25.

 \begin{itemize}
 \tightlist
 \item
 \emph{Summary:} An evaluation of Yau's systematic introduction of
 partial differential equations into global geometry and topology.
 \end{itemize}
\item
 Schoen, R. (2020). ``The contribution of Shing-Tung Yau to the
 positive mass theorem.'' \emph{Journal of Geometric Analysis}, 30(2),
 1200--1215.

 \begin{itemize}
 \tightlist
 \item
 \emph{Summary:} An assessment of Yau's impact on general relativity
 and the mathematical formulation of isolated systems in physics.
 \end{itemize}
\item
 Chern, S. S. (1984). ``The mathematical influence of Shing-Tung Yau.''
 \emph{Chinese Mathematical Society Publications}, 4, 1--5.

 \begin{itemize}
 \tightlist
 \item
 \emph{Summary:} Chern's reflections on the academic trajectory of
 his well-known student and his impact on geometric education and
 institution building.
 \end{itemize}
\item
 Lawson, H. B. (1983). ``Shing-Tung Yau: Fields Medalist.''
 \emph{Science}, 221(4616), 1151--1153.

 \begin{itemize}
 \tightlist
 \item
 \emph{Summary:} A general-audience overview of Yau's mathematical
 path and his unification of analysis, geometry, and physics.
 \end{itemize}
\item
 Witten, E. (2019). ``Shing-Tung Yau and string theory.'' \emph{Asian
 Journal of Mathematics}, 23(3), 353--360.

 \begin{itemize}
 \tightlist
 \item
 \emph{Summary:} Edward Witten's perspective on how Calabi-Yau
 geometry transformed mathematical physics and string
 compactification.
 \end{itemize}
\item
 Hamilton, R. S. (2000). ``Shing-Tung Yau and the development of
 geometric flows.'' \emph{International Press Publications}, 7, 25--40.

 \begin{itemize}
 \tightlist
 \item
 \emph{Summary:} An overview of Yau's influence on the formulation
 and analysis of the parabolic flows of Riemannian metrics.
 \end{itemize}
\item
 Uhlenbeck, K. (2010). ``The impact of Shing-Tung Yau's work on gauge
 theory.'' \emph{Journal of Differential Geometry}, 85(1), 1--15.

 \begin{itemize}
 \tightlist
 \item
 \emph{Summary:} An examination of the relationship between Hermitian
 Yang-Mills connections and stability on algebraic manifolds.
 \end{itemize}
\end{itemize}

\begin{center}\rule{0.5\linewidth}{0.5pt}\end{center}


\begin{center}\rule{0.5\linewidth}{0.5pt}\end{center}

\section*{See Also}
\addcontentsline{toc}{section}{See Also}

This chapter relates to the following work by other Fields Medalists:
\begin{itemize}
  \item \textbf{\hyperref[ch2006terencetao:terence-tao-fields-medal-2006-a-comprehensive-report-on-his-work-and-impact]{See Chapter on related work}} --- Tao worked on related PDEs.
  \item \textbf{\hyperref[ch1954kunihikokodaira:kunihiko-kodaira-fields-medal-1954-a-comprehensive-report]{See Chapter on related work}} --- Kodaira's work on harmonic integrals relates to Yau.
\end{itemize}

\emph{(End of report)}
